% Figure blocks. Each guards its own graphic so the manuscript still compiles while a
% figure is being regenerated; check_refs.py fails the build if any guard is still active
% at submission time.
\newcommand{\reviewfig}[4]{%
  \begin{figure}[htbp]\centering
  \IfFileExists{figures/#1}{\includegraphics[width=\textwidth]{figures/#1}}%
    {\fbox{\parbox[c][3cm][c]{0.9\textwidth}{\centering FIGURE PENDING: \texttt{\detokenize{#1}}}}}
  \caption{#2}\label{#3}
  \end{figure}#4}

\reviewfig{gb_fig1_atlas_overview.pdf}{\textbf{Sparse autoencoder feature atlas of Geneformer
V2-316M.} TopK sparse autoencoders trained on all 18 layers yield 82,525 alive features.
Reconstruction quality declines with depth while dead features increase.}{fig:overview}{}

\reviewfig{gb_fig4_ontology_heatmap.pdf}{\textbf{Biological annotation follows a U-shaped profile
across layers.} Early layers encode molecular programs, middle layers develop more abstract
representations, later layers re-specialize before becoming prediction-focused. Enrichment uses
Gene Ontology, KEGG (Kanehisa Laboratories), Reactome, STRING and TRRUST gene
sets.}{fig:ushape}{}

\reviewfig{gb_fig5_coactivation.pdf}{\textbf{Co-activation module structure across layers.}
(A)~Modules per layer. (B)~Fraction of alive features covered by a module. (C)~Co-activation edge
density.}{fig:modules}{}

\reviewfig{fig_highways.pdf}{\textbf{Cross-layer information flow is continuous.} (a)~Highway
fraction for the three long-range layer pairs after permutation-null correction; the
99th-percentile null threshold sits below the applied threshold for every pair, so the correction
does not shift the fractions. (b)~The same quantity for all 17 consecutive layer pairs, with no
sharp drop at any transition.}{fig:highways}{}

\reviewfig{fig_svd_capacity.pdf}{\textbf{The dictionary re-parameterizes the principal subspace
rather than reaching past it.} (a)~Cumulative projection $\rho_k$ of each decoder direction onto
the rank-$k$ principal subspace, median over features, against the random-direction expectation
$k/d$; the shaded band is the interquartile range at layer~17 and the dotted line marks the
50-axis basis. (b)~Truncated-SVD rank required to match the SAE's own variance explained, per
layer. (c)~Significant enrichment terms per direction when the leading 50 principal axes are
annotated with the protocol used for SAE features. Enrichment uses Gene Ontology, KEGG
(Kanehisa Laboratories) and Reactome gene sets.}{fig:svd}{}

\reviewfig{fig_concepts.pdf}{\textbf{Dictionary atoms and the distinct programs they express.}
(a)~Atoms and distinct gene-level programs per layer. (b)~Aggregate across all 18 dictionaries.
(c)~Mean absolute coherence between unit-normalized decoder atoms against the Welch bound, the
smallest achievable maximum coherence for that many unit vectors in 1,152
dimensions.}{fig:concepts}{}

\reviewfig{fig_causal.pdf}{\textbf{Causal feature ablation with an independent evaluation set.}
Distribution of specificity ratios across features, for three feature-selection rules and three
evaluation gene sets: the feature's top-20 genes intersected with its annotated term; the annotated
term's genes excluding that top-20; and a size-matched random gene set stratified on expression rank
and detection frequency, drawn per feature. The dotted line marks a ratio of one.}{fig:causal}{}

\reviewfig{fig_ceiling.pdf}{\textbf{What these data can support.} (a)~Fraction of transcription
factors whose predicted target set is enriched for its curated targets, at matched prediction-set
size and under one hypergeometric framework, for every method. (b)~How many perturbed
transcription factors have enough curated targets among the measured genes to be evaluable at
all.}{fig:ceiling}{}

\reviewfig{fig_power_replication.pdf}{\textbf{Unit of replication, sample size, and replication
across cell lines.} (a)~Responding features per target under a cell-level test against a
position-pooled test on the same cells; the dashed line is equality. (b)~Detection of the
knocked-down gene itself, and the number of responding features, as a function of perturbed cells
per target, restricted to the 14 factors with at least 100 cells so that sample size is not
confounded with which factors are available. (c)~Fraction of evaluable transcription factors with
target enrichment in each cell line, with perturbed and control cells drawn from the same
line.}{fig:power}{}
